\chapter{LITERATURE ANALYSIS}

Hyperspectral image (HSI) classification has been a focus area of remote sensing for many years, especially for being able to identify surface materials, something that broadband sensors are unable to do. There are challenges in the field that require a balance between deep models and the limited amount of training data that has been labeled. This section puts HSI classification methods into groups based on the models used and the challenges that they solve in order to help understand the development of these methods.

\section{Traditional Machine Learning Approaches}

In the beginning, HSI classification used shallow and statistical methods of machine learning. Dominant frameworks in this time were Support Vector Machines (SVM) and Random Forest (RF). An example of early work in HSI image classification was by Du et al. in using sparse discriminative representations and multinomial logistic regression to achieve 97.71% accuracy on the Indian Pines data set. In a similar work, Xia et al. used extended morphological profiles and ensemble techniques to achieve better results using RF. To solve the ensemble overfitting problem, a cascaded RF approach was proposed to organize multiple classifiers in consecutive layers. Numerous SVM variations were created, including probabilistic SVMs combined with guided filters, and spectral and spatial SVMs with Markov Random Field (MRF) approaches.

Although the methods mentioned previously are robust to noise and interpret output, they have become limited in terms of classification accuracy. A challenge for machine learning is to capture the associations in high-dimensional spectral spaces. Classical techniques in feature extraction also struggle to meet the accuracy of representations that deep networks learn.

\section{Convolutional Neural Network-Based Approaches}

Learning spatial and spectral features simultaneously and without manual feature engineering has changed HSI classification with CNNs. 3D-CNNs with L2 regularization and dropout were applied by Chen et al. and achieved a groundbreaking 99.66% classification accuracy for the Pavia University dataset, showing how well 3D convolutions work with spectral-spatial data cubes. In a different approach, 1D-CNNs were implemented by Li et al. to develop deep features of pixel pairs in order to mitigate the problem of having very few labeled data.

A challenge of 3D-CNNs is the parameter cost. Therefore, some strategies were introduced. AcousticSN was able to balance spectral-spatial expressiveness against computational cost by combining both 2D and 3D convolutions. In a different approach, Ahmad et al. introduced small 3D-CNN architectures with faster inference times that maintained the same accuracy. Boundary artifacts were addressed by Paoletti et al. by mirroring borders, which mitigated information loss for the edges of each patch. Combined, these approaches provided a strong justification for deep learning features with CNNs for HSI classification, but their requirements for labeled data and GPU memory grow with depth, which limits their use in real-world situations with a lack of data.

\section{Autoencoders, Generative Networks, and RNNs}

Data that is insufficiently labeled has prompted the use of architectures that give representations of a decent quality without much supervision. For instance, Stacked Autoencoders (SAE) offer a good method of performing unsupervised feature extraction by dimensionality reduction. Combining SAEs with CNNs, as done by Li et al., focuses on learning spatial weights that are adaptive to one locality, while Zhao et al. focused on SAE-based compression and 3D deep residual networks that combine both dimensionality reduction and classification, but assume a significant amount of labeled data.

Sequence modeling in HSI began with Recurrent Neural Networks. In their work, Shi et al. employed RNNs to incorporate spectral-spatial features of different scales. They considered different spectral bands to be in a temporal sequence. Another approach to the problem was Generative Adversarial Networks, which synthesized samples and improved small labeled datasets and, in turn, aimed to mitigate the generalization issue caused by the lack of labeled samples.

Each of these frameworks imposes additional issues. SAE and RNN frameworks suffer from the burden of sequential computation, and update the systems slowly and GANs offer unpredictable results. Small changes to hyperparameters can annihilate generators or result in mode collapse and create uninformative training data.

\section{Transformer Architectures and Semi-Supervised Methods}

In the realm of HSI classification, Transformers improve on capturing long range spectral-spatial relations. Their addition to CNN based local feature extraction has also produced positive results. Recent tokenization methods referring to the encoding of spectral-spatial areas as sequences have allowed more unmediated use of the Vision Transformer. Thanks to recent architectures featuring dual attention, benchmark results on various datasets have surpassed 99%.

Transformers have proven especially powerful in this field, but their increased architecture cost is not overshadowed by previous methods. Their cost and hardware demands are greater for large datasets and decrease their usefulness in low-data scenarios.

In the realm of semi-supervised learning, Sellami et al. showed that 3D deep networks, combined with adaptive band selection, leverage unlabeled data effectively and reduce the gap in accuracy of the labeled and unlabeled data. This research serves as an excellent proof of concept. Most of the problems this research identifies, like decision boundary bias regarding pseudo-label noise and confidence threshold selection, are still unresolved in the domain.

\section{Current Standing and Original Contribution of This Study}

A common theme in the reviewed research is that the models that are the most accurate when a lot of labeled data are available are also the most fragile models when labeled data are sparse, and vice versa. Transformers and deep CNNs tend to overfit when trained on small labeled datasets. Semi-supervised and generative approaches to the data problem introduced instability to the training process and sensitivity to thresholds, and no previous research has incorporated stability directly into the semi-supervised training process.

This research aims to solve this issue by combining CNNs and Transformers in a Self-Training framework incorporating Curriculum Learning. This is the technique that organizes the training samples in order of the estimated level of difficulty so that the model does not focus on the noisy pseudo-labels early on in the training process. In addition to its novelty, most of the other research in the domain consists of models that only function through command-line interfaces, which makes them inaccessible to most practitioners that do not have a specialized knowledge of deep learning. The work presented here aims to convert research models into applications that can be deployed and used outside of an academic environment through the introduction of a multithreaded PyQt6 GUI. It also aims to provide real-time predictions using other external datasets.
